\section{Unweighted joint index authentication and power pooling}
\label{sec:reciprocal23}
This theorem combines the incoming source 176 raw-index orientation argument
with the mixed-base two-integrality argument~\cite{joint176,mixed}. Neither ingredient is credited
as new here. The new synthesis removes the base weights from a joint
orientation and power certificate. It needs no common sign or pairwise
coprimality of the bases, and remains valid at zero and unit bases.
The emitted source instances separately authenticate their actual ports.

\subsection{The complete joint residual}
Let $W\ge258$, $|A|>W^W$, and let $b\ge16$ be an integer
with $12b+10\le W$.
Suppose $p\le15$ raw principal Pell pairs have classified representations
\[
 (X_i,Y_i)=\varepsilon_i(X_{\nu_i}(A),Y_{\nu_i}(A)),
 \qquad\varepsilon_i\in\{-1,1\},\quad |\nu_i|\le b.
\]
The requested signed indices obey $|k_i|\le b$. There are at most five
base blocks and at most 17 total requests. Every actual base $n_j$ and
proposed output $m_i$ has absolute value at most $b$. Requests may reuse a
principal pair. The retained gates give $Z_i=AY_i$. Define
\begin{align*}
 d&=A-1,\qquad T_j=2An_j-n_j^2-1,\\
 \Omega&=\sum_{i=1}^{p}\bigl[(X_i-1)^2+(Y_i-k_i)^2\bigr],\\
 H_j&=\sum_{i\text{ in block }j}Y_i(m_i-X_i),\\
 S_j&=\sum_{i\text{ in block }j}\bigl[(m_i-X_i+Z_i)^2-Y_i^2\bigr],
 \qquad E_j=S_j-2n_jH_j.
\end{align*}
The unweighted joint residual, with one new unrestricted integer $z$, is
\begin{equation}\label{eq:unweighted-joint}
 \Omega\prod_jT_j
 +d\sum_j E_j\prod_{h\ne j}T_h
 -zd\prod_jT_j=0.
\end{equation}
Products and sums here are fixed finite polynomial operations. The display
introduces no reciprocal gate or unspecified integer parameter.

\begin{theorem}\label{thm:unweighted-joint}
Under these hypotheses, an integer $z$ satisfying
\eqref{eq:unweighted-joint} exists if and only if every principal pair is
canonical at its requested signed index and every request at a nonzero
base satisfies $m_i=n_j^{k_i}$ in $\Q$. A base with $|n_j|\ge2$ forces $k_i\ge0$.
Unit bases allow signed indices. Zero-base requests impose no condition.
\end{theorem}
\begin{proof}
Put $M_0=W^W$. As $A$ is integral and $|A|>M_0$, $|d|\ge M_0$ and
$d$ has the sign of $A$. At $n_j=0$, $T_j=-1$. Otherwise
\[
 \frac{|T_j|}{|n_j|}
 \ge2|A|-\frac{n_j^2+1}{|n_j|}
 \ge2|A|-b-1>M_0,
\]
and $n_jT_j$ has the sign of $A$. In particular all denominators used
in this proof are nonzero.

Modulo $d$, the signed Pell recurrence gives $X_\nu\equiv1$ and
$Y_\nu\equiv\nu$. Therefore $\Omega\equiv o\pmod d$, where
\[
 o=\sum_i[(\varepsilon_i-1)^2+(\varepsilon_i\nu_i-k_i)^2],
 \qquad0\le o\le60(1+b^2)\le120b^2.
\]
The indices here are bounded \emph{classified} indices, not the requested
ones whose authenticity remains to be proved.

For each request put $\delta_i=m_i-X_i+(A-n_j)Y_i$. Exact expansion gives
\begin{equation}\label{eq:joint-error-identity}
 \sum_{i\text{ in }j}\delta_i^2=E_j-T_j\sum_{i\text{ in }j}Y_i^2.
\end{equation}
The Pell recurrence, evaluated at the square root $n_j-A$ modulo $T_j$,
gives for nonzero $n_j$
\[
 X_i+(n_j-A)Y_i\equiv\varepsilon_i n_j^{\nu_i}\pmod{T_j}.
\]
Negative powers exist in this residue ring because $\gcd(n_j,T_j)=1$.

\textbf{First integrality step.}
Divide~\eqref{eq:unweighted-joint} by $d\prod_jT_j$ and use
\eqref{eq:joint-error-identity}. It follows that
$\Omega/d+\sum_j\sum_i\delta_i^2/T_j$ is integral.
Remove the zero-base blocks; those summands are already integers.
Let $J$ be the remaining blocks and define
\[
 C_*=(\prod_{j\in J}n_j)^b,
 \qquad L_*=\prod_{j\in J}n_j,
\]
with empty products one. For a request in $J$, define the integer
\[
 e_i=C_*m_i-\varepsilon_i(C_*/n_j^b)n_j^{b+\nu_i},
 \qquad D_j=\sum_{i\text{ in }j}e_i^2,\quad D=\sum_{j\in J}D_j.
\]
All exponents are nonnegative. The recurrence congruence makes
$C_*\delta_i\equiv e_i\pmod{T_j}$. Hence
\[
 R_*=C_*^2o/d+\sum_{j\in J}D_j/T_j
\]
is an integer. Since $|C_*|\le b^{5b}$ and $|e_i|\le2b^{6b}$,
\begin{align*}
 C_*^2o+D
 &\le120b^{10b+2}+68b^{12b}\le188b^{12b+2}\\
 &<b^{12b+5}<W^{12b+5}\le W^W=M_0.
\end{align*}
Here $188<b^3$, $b<W$, and $12b+5\le W$ by hypothesis.
Each active denominator has magnitude at least $M_0$, so $|R_*|<1$.
Thus $R_*=0$, although its summands may have different signs.

\textbf{Second integrality step.}
The two exact rational identities
\[
 \frac{2A}{d}=2+\frac2d,
 \qquad\frac{2A}{T_j}=\frac1{n_j}+
             \frac{n_j^2+1}{n_jT_j}
\]
split $2AR_*=0$ as $U_*+V_*=0$, where
\begin{align*}
 U_*&=2C_*^2o+\sum_{j\in J}D_j/n_j,\\
 V_*&=2C_*^2o/d+
       \sum_{j\in J}\frac{D_j(n_j^2+1)}{n_jT_j}.
\end{align*}
The integer $L_*U_*$ makes $L_*V_*=-L_*U_*$ integral. The denominator
estimate also gives
\[
 \frac{n_j^2+1}{|n_jT_j|}
 \le\frac{1+n_j^{-2}}{2|A|-b-1}<\frac2{M_0}.
\]
Using $|L_*|\le b^5$ yields
\begin{align*}
 |L_*V_*|
 &\le\frac{240b^{10b+7}+136b^{12b+5}}{M_0}
 \le\frac{376b^{12b+7}}{M_0}\\
 &<\frac{b^{12b+10}}{M_0}
 <\frac{W^{12b+10}}{M_0}\le1.
\end{align*}
We used $376<b^3$, $b<W$, and $12b+10\le W$.
Thus $L_*V_*=0$; since $L_*\ne0$, $V_*=0$.
Every nonzero summand of $V_*$ has the sign of $A$: $o,D_j\ge0$,
$d$ has that sign, and so does $n_jT_j$. Consequently $o=0$ and every
$D_j=0$. The first forces $\varepsilon_i=1$ and $\nu_i=k_i$ for every
principal pair. The second gives $e_i=0$, hence $m_i=n_j^{k_i}$.
For $|n_j|\ge2$, negative $k_i$ would force the impossible integer equation
$m_in_j^{-k_i}=1$. If $J$ is empty, the same proof has $C_*=L_*=1$
and $V_*=2o/d$, so it still authenticates every pair.

Conversely, canonical pairs give $d\mid\Omega$. True nonzero-base powers
make every $\delta_i$ divisible by $T_j$, so~\eqref{eq:joint-error-identity}
gives $T_j\mid E_j$. At zero bases this is automatic because $T_j=-1$.
Thus $z=\Omega/d+\sum_jE_j/T_j$ is integral and satisfies the residual.
\end{proof}

\subsection{Square and affine frame corollaries}
For the exact square frame, $b=\lfloor\sqrt{W-1}\rfloor$ has
$b\ge16$ and $12b+10\le b^2+1\le W$: the first difference is
$b^2-12b-9$, positive at 16 and increasing thereafter. The original
square-frame interface is therefore a special case.

A centered frame could instead supply $b\le14\sqrt W$ and
$W\ge2^{42}$. Then $12b+10\le168\sqrt W+10\le W$.
This is a conditional extension: every actual base, output, nominal index
and raw classified index must be independently bounded at that same $b$.
It is not an automatic port transfer to an unaudited weak frame. The selected
reciprocal endpoints retain the conservative exact-square source.

\subsection{Restore the source before recovering its signs}
The source's orientation quotient is $\Omega/d$, with any explicitly recorded
retained-row correction restored. For a singleton whose rows are
$T=T_j$ and $\delta_i+TQ=0$, restore $T=T_j$, $Q=-\delta_i/T_j$.
For an aggregate restore
\[
 L_j=E_j/T_j,\qquad K_j=2H_j+(2A-n_j)L_j.
\]
These are integers also at zero or unit bases. Substitution in the exact
identities restores every original row. Only then apply the original
semantic bootstrap, which proves the actual arithmetic bases nonzero and
positive in its own established order. No circular assumption that $q$ is
already positive is needed for the joint theorem.

One may partition the base blocks and fuse orientation into one part using
\eqref{eq:unweighted-joint}. That part authenticates \emph{all} principal
pairs, even if it contains only some requests. The other parts then use the
incoming unweighted mixed-base theorem at already canonical pairs. Its
proof is the two integrality steps above with $o=0$, so no new sign or
coprimality premise is introduced. All quotient restorations are performed
before source decoding.

\subsection{Authenticating actual ports and the retained frame}
The weighted reference has 14 principal pairs, three singleton requests and
two aggregate blocks. Its certificate checks every norm and multiplication
gate, all private coordinates, and the exact orientation identity
\[
 \Omega-(R_{\rm orient}+(A-1)q_{\rm orient})
       =\sum_{i=1}^{14}R_{{\rm norm},i}.
\]
The original squared $C$ cap bounds every raw classified index by $|C|\le b$
\emph{before} orientation. Its frame contains the actual growth index and
all requested bases, indices and outputs. Its inherited $q\equiv1\pmod3$
restriction belongs to that frozen source's completeness construction, not
to Theorem~\ref{thm:unweighted-joint} itself.

The new affine-source instance instead binds the conservative $m=16$
209-coordinate source of Section~\ref{sec:quartic23}, with 13 principal
pairs and four remaining power blocks. It checks the actual affine ports
and all retained-row corrections from the centered norm/orientation chart.
This is the same theorem with fewer active pairs, not an assumption that
similar coordinate names denote equal polynomials. In particular, the
weak affine-atom frame discussed there is \emph{not} substituted into the
square-frame hypothesis of the joint theorem.

\subsection{A retained row changes the numerator chart}
The smallest selected endpoint uses a further exact algebraic adjustment,
not a weaker-frame theorem. Let $R_W=0$ be the retained frame residual of
the conservative source and replace the orientation numerator by
\[
 \Omega'=\Omega+R_W.
\]
All retained source zeros have $R_W=0$, so the joint zero-set argument is
unchanged. In the actual centered source coordinates this addition cancels
the quadratic $x^2$ term. A later monic substitution for $x$ can then have
degree four without making the selected pool row exceed degree eight.
The certificate checks the exact coefficient cancellation and restores the
retained row before invoking Theorem~\ref{thm:unweighted-joint}.
It is not a claim that $\Omega'$ is itself nonnegative away from those
retained zeros.

For the 113-variable endpoint the four base blocks are partitioned as
$(B)$ and $(D,q,q+1)$, with orientation fused into the first part. The
source has thirteen principal pairs and fifteen power requests, not the
fourteen-pair weighted reference. The 202-coordinate, 142-row joint seed
then undergoes 89 exact unit-pivot eliminations, leaving 113 coordinates
and 53 residuals of degree at most eight. Its sum of squares has exact
degree 16, support 1030202 and height $16250018885420533678080$.

The sparse 114-variable companion keeps orientation separate and partitions
the power blocks as $(B,D)$ and $(q,q+1)$. It has 54 residuals, exact degree
16, 127818 terms and the same height. This is a useful sparseness tradeoff,
not another use of the fused orientation theorem. The 138-variable companion
has 78 residuals, exact degree eight, 12099 terms and height $2^{50}$.
The 138-variable point uses direct monic elimination from the conservative
source, without a new pool, and is not credited to the joint theorem.
All three full collected polynomials are supplied; their supports are not
estimates inferred from circuit sizes.

\subsection{The finite substitution specification}
For clarity, a unit-pivot step at the current residual row has the form
\[
 Q_r=\epsilon x_v+f(\widehat{x_v}),\qquad\epsilon\in\{-1,1\},
\]
with no other occurrence of $x_v$ in $Q_r$. Substitute $x_v=-\epsilon f$
in every other residual, then delete row $r$ and coordinate $v$. This is
an integral zero-set bijection: insertion of that polynomial is its exact
inverse, and the deleted equation becomes an identity. Later restorations
compose in reverse order. Every step's resulting degree is checked after
full exact coefficient collection; a generic degree-propagation bound is
not used as an attained degree certificate.

The appendix gives the complete conservative source, literal principal-pair
and request ports, exact removed rows, fixed pool partitions, and each
ordered triple $(v,r,\epsilon)$. Starting with the stated seed, set
$\sigma_0$ to the identity and successively apply
\[
 \sigma_{j+1}=\sigma_j
 \bigl[x_v\mapsto-\epsilon\,\sigma_j(Q_r)|_{x_v=0}\bigr].
\]
Here the substitutions mean applying the current pivot replacement to every
current row; the recorded order removes any ambiguity about composition.
The final polynomial is the sum of squares of all retained substituted
rows, in the explicitly listed retained coordinate order. This is a finite
ordinary polynomial definition with no existential register or division.
The supplied full residual and coefficient arrays independently agree with
this specification.

The matched unfused control uses the same source and admissible monic pivot
strategy. It isolates the extra elimination enabled by the joint numerator
chart. The sparse 114-variable endpoint and the low-degree 138-variable
endpoint are separate controls with their own partitions and metrics; they
are not falsely reported as a simultaneous improvement in all objectives.
The unbounded joint theorem, exact substitutions and inherited source
calibration together establish~\eqref{eq:calibration} for every selected
endpoint on $\Z$.

